\documentclass[11pt,a4paper]{article}

\usepackage{amsmath,amssymb}
\usepackage{booktabs}
\usepackage{array}
\usepackage{longtable}
\usepackage{graphicx}
\usepackage{float}
\usepackage{caption}
\usepackage{subcaption}
\usepackage[margin=1in]{geometry}
\usepackage{xcolor}
\usepackage{hyperref}
\usepackage{enumitem}
\usepackage{fancyhdr}
\usepackage{titlesec}
\usepackage{tabularx}
\usepackage{pdflscape}
\usepackage{microtype}
\usepackage{setspace}

\hypersetup{colorlinks=true,linkcolor=black,urlcolor=blue,citecolor=black}
\setstretch{1.08}
\setlength{\parindent}{0pt}
\setlength{\parskip}{6pt}
\renewcommand{\arraystretch}{1.15}

\pagestyle{fancy}
\fancyhf{}
\lhead{Missing Data \& Statistical Inference}
\rhead{Project Report}
\cfoot{\thepage}
\setlength{\headheight}{14pt}

\titleformat{\section}{\Large\bfseries}{\thesection.}{0.6em}{}
\titleformat{\subsection}{\large\bfseries}{\thesubsection}{0.6em}{}
\titleformat{\subsubsection}{\normalsize\bfseries}{\thesubsubsection}{0.6em}{}

\begin{document}

% ---------------- COVER ----------------
\begin{titlepage}
\thispagestyle{empty}
\centering
\vspace*{1.25in}
{\Huge\bfseries MISSING DATA \&\\ STATISTICAL INFERENCE\par}
\vspace{0.35in}
{\Large A Simulation-Based Study of MCAR, MAR and MNAR\par}
\vspace{0.25in}
{\large Comparing Seven Missing-Data Handling Methods Within Each Mechanism\par}
\vspace{1.0in}
{\Large \textbf{Project Report}\par}
\vspace{0.2in}
{\large M.Sc. Statistics\par}
\vfill
{\large September--October 2026\par}
\end{titlepage}

% ---------------- ABSTRACT ----------------
\pagenumbering{roman}
\section*{Abstract}
This project investigates how missing-data handling methods affect regression estimation and statistical inference when observations are incomplete. A complete dataset of $n=100$ observations was simulated from a known regression model,
\[
Y = 5 + 2X_1 - 1.5X_2 + X_3 + \epsilon,
\]
so that the true coefficients were known. Missing values were then deliberately introduced in the response variable $Y$ under three missingness mechanisms: Missing Completely At Random (MCAR), Missing At Random (MAR), and Missing Not At Random (MNAR).

For each mechanism, the same incomplete dataset was given to seven methods: complete-case analysis, mean imputation, hot-deck imputation, regression imputation, stochastic regression imputation, Multiple Imputation by Chained Equations (MICE) using predictive mean matching (PMM), and the Expectation-Maximization (EM) algorithm. The central comparison is therefore \textbf{within each missingness mechanism}: MCAR methods are compared with one another, MAR methods with one another, and MNAR methods with one another.

The regression model estimates all four coefficients, but the detailed inferential analysis focuses on $\beta_1$, the coefficient of $X_1$, whose true value is $2$. Performance is examined using the estimated coefficient, bias, bootstrap standard error, and 95\% confidence interval. MICE additionally combines uncertainty across five completed datasets using Rubin's rules.

The clearest pattern in the simulation is the large downward distortion produced by mean imputation. Under MCAR, for example, mean imputation produced $\hat{\beta}_1=1.089$, compared with the true value of $2$, whereas complete-case analysis produced $1.785$. Under MAR, MICE produced a pooled estimate of $1.797$, while mean imputation produced $1.307$. Under MNAR, all methods remained below the true value in this particular realization, with mean imputation again producing the largest deviation. These results demonstrate that missing-data handling is part of the statistical analysis itself rather than merely a data-cleaning step.

\newpage
\tableofcontents
\newpage
\pagenumbering{arabic}

% ---------------- 1 ----------------
\section{Introduction}
Missing data are a common feature of real statistical datasets. Once observations are incomplete, the analysis may change not only because fewer observations are available, but also because the way missing observations are handled can alter the relationships among variables and the uncertainty of the resulting estimates.

A central distinction in missing-data analysis is the mechanism that generated the missingness. The standard classification is MCAR, MAR, and MNAR. Under MCAR, missingness does not depend on the data. Under MAR, missingness may depend on information that remains observed. Under MNAR, missingness may also depend on the value that is itself missing. The course material presents these mechanisms as a nested hierarchy,
\[
\text{MCAR}\subset\text{MAR}\subset\text{MNAR},
\]
with MCAR the most restrictive condition and MNAR the most general. 

The practical issue is that there is no single way to handle missing observations. A researcher might delete incomplete cases, replace missing values with a mean, borrow values from donors, predict them from a regression model, create multiple imputations, or estimate model parameters directly using likelihood-based methods. These choices can lead to different estimates and different inferential conclusions.

The purpose of this project is therefore not to treat MCAR, MAR and MNAR as three competing ``methods.'' They are treated as three experimental conditions. The main comparison is between the seven missing-data handling approaches \textbf{within each condition}. This gives the project a clear structure: hold the missingness mechanism fixed, give every method the same incomplete dataset, and investigate how the choice of method changes the analysis.

% ---------------- 2 ----------------
\section{Objective and Research Question}
\subsection{Objective}
The objective is to examine how different missing-data handling methods affect regression estimation and statistical inference when applied to the same incomplete dataset under a specified missingness mechanism.

The study evaluates both the location of the estimate and the uncertainty around it. In particular, the analysis examines:
\[
\hat{\beta}_1,
\qquad
\text{Bias}(\hat{\beta}_1),
\qquad
SE(\hat{\beta}_1),
\qquad
95\%\;CI.
\]

Although the fitted regression contains four coefficients, $\beta_1$ is selected as the primary parameter of interest for the detailed inference comparison. Its true value is known to be
\[
\boxed{\beta_1=2}.
\]

\subsection{Research Question}
\textbf{How does the missing-data handling method affect the estimated regression relationship and the associated statistical uncertainty when the missingness mechanism is held fixed?}

This question is examined three times: once for MCAR, once for MAR, and once for MNAR.

% ---------------- 3 ----------------
\section{Data-Generating Model}
To create a controlled experiment, a complete dataset was generated before any values were made missing. The response was generated from the linear model
\[
Y_i = \beta_0 + \beta_1X_{1i} + \beta_2X_{2i} + \beta_3X_{3i} + \epsilon_i.
\]

The parameters used in the simulation were
\[
\beta_0=5,\qquad
\beta_1=2,\qquad
\beta_2=-1.5,\qquad
\beta_3=1.
\]

The predictors and random error were generated as
\[
X_1,X_2,X_3\overset{ind}{\sim}N(0,1),
\qquad
\epsilon\sim N(0,1).
\]

A sample of $n=100$ observations was generated. The complete dataset therefore represents the hypothetical data that would have been available if no information were missing.

The benefit of simulation is that the true parameter is known. With real data, the true value of a regression coefficient is not available as a benchmark. Here, however, an estimate such as $\hat\beta_1=1.70$ can be directly compared with the known truth $\beta_1=2$.

\subsection{Baseline Complete-Data Regression}
Before introducing missingness, the regression was fitted to the complete dataset. The resulting coefficients were:

\begin{table}[H]
\centering
\caption{Baseline complete-data regression.}
\begin{tabular}{lrr}
\toprule
Parameter & True value & Complete-data estimate\\
\midrule
$\beta_0$ & 5.000 & 5.0875\\
$\beta_1$ & 2.000 & 1.8227\\
$\beta_2$ & -1.500 & -1.5382\\
$\beta_3$ & 1.000 & 1.0269\\
\bottomrule
\end{tabular}
\end{table}

The fact that the complete-data estimate of $\beta_1$ is $1.8227$ rather than exactly $2$ is expected. The simulated dataset is a sample, so sampling variation exists even before missingness is introduced. This is important when interpreting later deviations from the true value: not all deviation from $2$ is necessarily caused by missing data.

% ---------------- 4 ----------------
\section{Missing-Data Mechanisms}
After generating the complete dataset, missingness was introduced only in $Y$. The same complete dataset was used as the starting point for all three mechanisms. This means that, within each mechanism, all seven methods receive the same incomplete data.

\subsection{MCAR}
Under MCAR, every observation has the same probability of having its $Y$ value made missing. The simulation used
\[
P(R_i=0)=0.20,
\]
where $R_i=0$ denotes a missing response.

For each of the 100 observations, a random number between 0 and 1 was generated. Whenever the number was below $0.20$, the corresponding $Y$ value was replaced by `NaN`.

In this particular realization, 25 observations had missing $Y$, so 75 observations retained an observed response. The realized percentage was therefore 25\%, even though the specified probability was 20\%. A probability of 20\% does not guarantee exactly 20 missing observations in a random sample of 100.

\subsection{MAR}
Under MAR, missingness depends on observed information. The simulation made the missingness probability depend on $X_1$:
\[
P(R_i=0\mid X_{1i})=
\begin{cases}
0.30,&X_{1i}<0,\\
0.10,&X_{1i}\geq 0.
\end{cases}
\]

Thus, observations with negative $X_1$ values were more likely to lose their $Y$ value than observations with non-negative $X_1$.

In this realization, 25 observations had missing $Y$. A useful diagnostic confirmed the intended relationship: the mean $X_1$ among observations with observed $Y$ was approximately $0.055$, while the mean $X_1$ among observations with missing $Y$ was approximately $-0.581$.

\subsection{MNAR}
Under MNAR, missingness depends directly on the value that becomes missing. The simulation first standardized the original response:
\[
Y_i^*=\frac{Y_i-\bar Y}{s_Y}.
\]

The probability of missingness was then defined as
\[
p_i=\operatorname{expit}\left(-1.3863+0.5Y_i^*\right),
\]
where
\[
\operatorname{expit}(z)=\frac{1}{1+e^{-z}}.
\]

Because $Y_i^*$ appears directly in the probability formula, observations with larger original $Y$ values have a higher probability of becoming missing. The probabilities are then compared with a new random number for each observation.

In the particular realization, 16 observations had missing $Y$. Using the original complete dataset, the mean $Y$ among observations that became missing was approximately $5.157$, while the mean among observations that remained observed was approximately $4.888$. This is consistent with the direction built into the MNAR mechanism.

\begin{table}[H]
\centering
\caption{Missingness mechanisms and realized sample sizes.}
\begin{tabular}{lccc}
\toprule
Mechanism & Dependence & Missing $Y$ & Observed $Y$\\
\midrule
MCAR & None & 25 & 75\\
MAR & Observed $X_1$ & 25 & 75\\
MNAR & Original $Y$ & 16 & 84\\
\bottomrule
\end{tabular}
\end{table}

The key distinction is therefore not the percentage of missing observations. The mechanism is defined by \textit{what the probability of missingness depends on}.

% ---------------- 5 ----------------
\section{Missing-Data Handling Methods}
For each missingness mechanism, the same seven approaches were applied.

\subsection{Complete-Case Analysis}
Rows with missing $Y$ were removed and the regression was fitted only to the remaining observations. No artificial values were created, but information was discarded.

Under MCAR, complete-case observations form a random subset of the complete data in the usual sense. Under MAR or MNAR, the validity of complete-case analysis requires stronger assumptions; in general, the course material emphasizes MCAR for the usual unbiasedness result.

\subsection{Mean Imputation}
Each missing response was replaced by the mean of the observed responses:
\[
Y_i^{imp}=\bar Y_{obs}
\quad\text{for missing }Y_i.
\]

This is simple, but all missing observations receive the same value. The course material shows that under MCAR this reduces the variance of the imputed outcome and attenuates correlations with predictors. Consequently, mean imputation can weaken regression relationships.

\subsection{Hot-Deck Imputation}
For each missing observation, an observed donor with similar values of $(X_1,X_2,X_3)$ was found using nearest-neighbour squared distance,
\[
D_{ij}=(X_{1i}-X_{1j})^2+(X_{2i}-X_{2j})^2+(X_{3i}-X_{3j})^2.
\]
The $Y$ value of the closest observed donor was then copied to the missing observation.

Because donor values are actual observed outcomes, hot-deck does not replace every missing value with the same constant. However, it remains a single-imputation method and therefore does not fully represent imputation uncertainty.

\subsection{Regression Imputation}
A regression model was fitted to observations with observed $Y$ values:
\[
Y=\beta_0+\beta_1X_1+\beta_2X_2+\beta_3X_3+\epsilon.
\]
For a missing response, the fitted model supplied
\[
\hat Y_i=\hat\beta_0+\hat\beta_1X_{1i}+\hat\beta_2X_{2i}+\hat\beta_3X_{3i}.
\]
The predicted value was then used as the imputed outcome.

\subsection{Stochastic Regression Imputation}
Stochastic regression imputation added residual variation to the regression prediction:
\[
Y_i^{imp}=\hat Y_i+e_i^*,
\]
where $e_i^*$ was sampled from observed regression residuals. This allows the imputed observations to retain random outcome variation rather than lying exactly on the fitted regression relationship.

\subsection{MICE Using Predictive Mean Matching}
MICE normally handles several incomplete variables by cycling through them one at a time, fitting a suitable conditional model at each step. In this project only $Y$ is incomplete, so the implementation is a simplified PMM-based multiple-imputation procedure.

The project used
\[
m=5
\]
completed datasets and
\[
k=5
\]
nearest potential donors in predictive mean matching. Five cycles were used in the current notebook implementation. Because only $Y$ is incomplete and the predictors are fully observed, this should be understood as a simplified one-variable MICE/PMM setting rather than a general multi-variable MICE workflow.

The initial missing values were filled with a temporary mean. A regression of observed $Y$ on the fully observed predictors was then fitted. For each missing observation, the predicted value was compared with predicted values of observed cases; the five closest observed cases were treated as possible donors, and one donor's actual observed $Y$ was randomly selected. This process generated one completed dataset. The process was repeated to generate five completed datasets.

Each completed dataset was analyzed separately. The five resulting $\hat\beta_1$ values were then pooled using Rubin's rules. If $W$ denotes the average within-imputation variance and $B$ the between-imputation variance, the total variance is
\[
T=W+\left(1+\frac{1}{m}\right)B.
\]
The pooled standard error is
\[
SE_{pooled}=\sqrt{T}.
\]

\subsection{Expectation-Maximization}
EM estimates the parameters directly from the incomplete data. The E-step calculates expected quantities involving the missing responses using current parameter estimates, and the M-step updates the regression coefficients and error variance.

The algorithm was iterated until changes in the parameters and variance were sufficiently small. In this project the final EM results were followed by bootstrap inference for $\beta_1$.

% ---------------- 6 ----------------
\section{Evaluation Framework}
\subsection{Bias}
The bias in a single simulation is measured relative to the known truth:
\[
\text{Bias}(\hat\beta_1)=\hat\beta_1-\beta_1=\hat\beta_1-2.
\]
Absolute bias is
\[
|\hat\beta_1-2|.
\]
A negative bias means that the estimate is below 2.

\subsection{Bootstrap Standard Error}
For the bootstrap-based methods, 500 bootstrap samples were generated by resampling observations with replacement within the corresponding method-specific workflow. The analysis was repeated for each bootstrap sample, giving
\[
\hat\beta_1^{(1)},\hat\beta_1^{(2)},\ldots,\hat\beta_1^{(500)}.
\]
The bootstrap standard error was estimated as
\[
SE_{boot}=SD\left(\hat\beta_1^{(1)},\ldots,\hat\beta_1^{(500)}\right).
\]

\subsection{95\% Confidence Intervals}
The bootstrap confidence intervals use the percentile method:
\[
CI_{95\%}=\left(q_{0.025},q_{0.975}\right),
\]
where $q_p$ denotes the $p$th percentile of the bootstrap distribution.

For MICE, the interval is obtained from the pooled estimate and Rubin's rules rather than treating the five imputations as one ordinary bootstrap sample.

\subsection{Interpretation of a Single Interval}
Because the project uses one simulated dataset for each mechanism, whether one confidence interval contains $\beta_1=2$ is only a statement about that particular interval. It is not an empirical coverage estimate. Long-run coverage would require repeatedly generating new complete datasets and repeating the entire missingness-and-analysis process.

% ---------------- 7 ----------------
\section{Results}
The results are organized by missingness mechanism. Within each mechanism, all seven handling methods are compared on the same incomplete dataset.

\subsection{MCAR: Compare All Seven Methods}
The MCAR dataset contained 25 missing $Y$ values and 75 observed responses. The results for the focal coefficient are shown below.

\begin{table}[H]
\centering
\small
\caption{MCAR: comparison of all seven methods.}
\begin{tabular}{lrrrrr}
\toprule
Method & $\hat\beta_1$ & SE & CI lower & CI upper & Bias\\
\midrule
Complete Case & 1.785 & 0.136 & 1.527 & 2.025 & -0.215\\
Mean Imputation & 1.089 & 0.229 & 0.649 & 1.565 & -0.911\\
Hot-Deck & 1.658 & 0.165 & 1.262 & 1.918 & -0.342\\
Regression Imputation & 1.785 & 0.136 & 1.527 & 2.025 & -0.215\\
Stochastic Regression & 1.706 & 0.140 & 1.487 & 2.046 & -0.294\\
MICE & 1.659 & 0.114 & 1.433 & 1.885 & -0.341\\
EM & 1.785 & 0.136 & 1.527 & 2.025 & -0.215\\
\bottomrule
\end{tabular}
\end{table}

Mean imputation is the clear outlier in this particular MCAR analysis. Its estimate is $1.089$, a difference of about $0.91$ from the truth. Its confidence interval also remains far below 2.

Complete-case analysis gives $1.785$, only about $0.038$ below the complete-data estimate of $1.823$. Regression imputation and EM give essentially the same coefficient in this particular setup. Hot-deck and MICE both give estimates around $1.66$, while stochastic regression gives $1.706$.

An important inferential difference is visible in the intervals. Complete-case, regression-imputation, stochastic-regression and EM intervals extend to or beyond 2, while the mean, hot-deck and MICE intervals stop below 2 in this particular run.

\begin{figure}[H]
\centering
\includegraphics[width=0.87\textwidth]{figures/ci_mcar.png}
\caption{MCAR: $\beta_1$ estimates and 95\% confidence intervals by method.}
\end{figure}

\subsection{MAR: Compare All Seven Methods}
The MAR dataset also contained 25 missing $Y$ values. Here, missingness was related to the observed predictor $X_1$.

\begin{table}[H]
\centering
\small
\caption{MAR: comparison of all seven methods.}
\begin{tabular}{lrrrrr}
\toprule
Method & $\hat\beta_1$ & SE & CI lower & CI upper & Bias\\
\midrule
Complete Case & 1.777 & 0.136 & 1.532 & 2.055 & -0.223\\
Mean Imputation & 1.307 & 0.186 & 0.915 & 1.659 & -0.693\\
Hot-Deck & 1.787 & 0.148 & 1.416 & 1.996 & -0.213\\
Regression Imputation & 1.777 & 0.136 & 1.532 & 2.055 & -0.223\\
Stochastic Regression & 1.759 & 0.147 & 1.499 & 2.065 & -0.241\\
MICE & 1.797 & 0.094 & 1.612 & 1.982 & -0.203\\
EM & 1.777 & 0.136 & 1.532 & 2.055 & -0.223\\
\bottomrule
\end{tabular}
\end{table}

Mean imputation again gives a substantially lower estimate, $1.307$, with bias $-0.693$. Its interval, $(0.915,1.659)$, lies entirely below 2.

The other methods are clustered much more closely. The MICE pooled estimate is $1.797$, the closest to 2 among the seven methods in this particular MAR realization, with absolute bias $0.203$. Hot-deck gives $1.787$, stochastic regression $1.759$, and complete case, regression imputation and EM all give approximately $1.777$.

For MICE, the five imputed datasets produced
\[
1.7683,\quad 1.8304,\quad 1.7944,\quad 1.7991,\quad 1.7925,
\]
which pool to
\[
\hat\beta_1=1.7969.
\]
The within-imputation variance was $0.008294$ and the between-imputation variance was $0.000493$, producing a pooled SE of $0.0943$. The 95\% interval was $(1.612,1.982)$.

\begin{figure}[H]
\centering
\includegraphics[width=0.87\textwidth]{figures/ci_mar.png}
\caption{MAR: $\beta_1$ estimates and 95\% confidence intervals by method.}
\end{figure}

\subsection{MNAR: Compare All Seven Methods}
The MNAR dataset contained 16 missing $Y$ values. Missingness was directly related to the original response value.

\begin{table}[H]
\centering
\small
\caption{MNAR: comparison of all seven methods.}
\begin{tabular}{lrrrrr}
\toprule
Method & $\hat\beta_1$ & SE & CI lower & CI upper & Bias\\
\midrule
Complete Case & 1.749 & 0.122 & 1.501 & 1.981 & -0.251\\
Mean Imputation & 1.413 & 0.170 & 1.047 & 1.702 & -0.587\\
Hot-Deck & 1.722 & 0.123 & 1.469 & 1.954 & -0.278\\
Regression Imputation & 1.749 & 0.122 & 1.501 & 1.981 & -0.251\\
Stochastic Regression & 1.719 & 0.124 & 1.517 & 1.987 & -0.281\\
MICE & 1.722 & 0.108 & 1.510 & 1.933 & -0.278\\
EM & 1.749 & 0.122 & 1.501 & 1.981 & -0.251\\
\bottomrule
\end{tabular}
\end{table}

Mean imputation again produces the largest deviation from 2, with $\hat\beta_1=1.413$ and bias $-0.587$. The other six methods lie between about 1.719 and 1.749.

The MICE results were
\[
1.7640,\quad1.6876,\quad1.7478,\quad1.7083,\quad1.7004,
\]
with pooled estimate
\[
\hat\beta_1=1.7216.
\]
The within-imputation variance was $0.010295$, the between-imputation variance was $0.001067$, and the pooled SE was $0.1076$. The resulting interval was $(1.510,1.933)$.

All seven MNAR intervals fall below 2 in this particular simulation. This is a notable empirical feature of the realized dataset, but it should not be interpreted as a general statement about the coverage of the methods.

\begin{figure}[H]
\centering
\includegraphics[width=0.87\textwidth]{figures/ci_mnar.png}
\caption{MNAR: $\beta_1$ estimates and 95\% confidence intervals by method.}
\end{figure}

\subsection{Cross-Method Summary}
The main conclusions are obtained by comparing methods within each mechanism.

\begin{table}[H]
\centering
\small
\caption{Absolute bias of $\beta_1$ within each missingness mechanism.}
\begin{tabular}{lrrr}
\toprule
Method & MCAR & MAR & MNAR\\
\midrule
Complete Case & 0.215 & 0.223 & 0.251\\
Mean Imputation & 0.911 & 0.693 & 0.587\\
Hot-Deck & 0.342 & 0.213 & 0.278\\
Regression Imputation & 0.215 & 0.223 & 0.251\\
Stochastic Regression & 0.294 & 0.241 & 0.281\\
MICE & 0.341 & 0.203 & 0.278\\
EM & 0.215 & 0.223 & 0.251\\
\bottomrule
\end{tabular}
\end{table}

This table should be read column by column. Within MCAR, for example, all seven entries describe methods applied to the same MCAR dataset. The same interpretation applies to MAR and MNAR.

The dominant pattern is the large absolute bias of mean imputation. In contrast, the other methods are more tightly clustered in all three conditions. The exact numerical ordering changes between mechanisms, which is another reason not to interpret this single simulation as a universal ranking.

\begin{figure}[H]
\centering
\includegraphics[width=0.82\textwidth]{figures/overview_beta1_grouped.png}
\caption{Overview of $\beta_1$ estimates, with MCAR, MAR and MNAR shown on the x-axis and the seven methods grouped within each mechanism.}
\end{figure}

\begin{figure}[H]
\centering
\includegraphics[width=0.72\textwidth]{figures/absolute_bias_heatmap.png}
\caption{Absolute bias of $\beta_1$. Lower values indicate estimates closer to the known true value of 2 in this particular simulation.}
\end{figure}

% ---------------- 8 ----------------
\section{Discussion}
\subsection{Mean Imputation Stands Out}
Across MCAR, MAR and MNAR, mean imputation is the method that most clearly changes the regression result. It replaces all missing responses with a common value, removing variation from the imputed part of the response distribution. In the simulation this produces a strong downward movement in $\hat\beta_1$.

Under MCAR, the estimate drops to $1.089$; under MAR it is $1.307$; under MNAR it is $1.413$. The absolute biases are therefore $0.911$, $0.693$ and $0.587$, respectively. These are substantially larger than the corresponding biases for the other methods in the same datasets.

The result is consistent with the theoretical course material, which shows that mean imputation shrinks outcome variance and attenuates correlations under MCAR. In a regression setting, reducing the natural variation in $Y$ can weaken the apparent relationship between $Y$ and its predictors.

\subsection{Complete Case, Regression Imputation and EM}
In this particular setup, complete-case analysis, regression imputation and EM produce nearly identical point estimates. This is not because the three procedures are the same. They use different ideas, but our model has a special structure: only $Y$ is incomplete and all predictors are fully observed.

Complete-case analysis simply discards the missing-response rows. Regression imputation predicts their $Y$ values from the observed predictors. EM works through expected quantities for the missing outcomes while iteratively updating the model. Despite these conceptual differences, their fitted coefficient estimates coincide in this simulation.

This is an important example of why numerical equality does not necessarily imply methodological equality.

\subsection{Hot-Deck and Stochastic Regression}
Hot-deck retains actual observed $Y$ values by borrowing values from similar donors. Stochastic regression also preserves additional outcome variability through sampled residuals. Both therefore behave differently from deterministic mean replacement.

The MCAR hot-deck estimate was $1.658$ and the stochastic-regression estimate was $1.706$. Under MAR they were $1.787$ and $1.759$, respectively. Under MNAR they were $1.722$ and $1.719$. The estimates are reasonably close to one another in several cases, but their confidence intervals and exact sampling behavior differ.

\subsection{MICE and Imputation Uncertainty}
MICE is conceptually important because it treats the missing values as uncertain rather than pretending there is one uniquely correct completed dataset. Five completed datasets were generated, producing five values of $\hat\beta_1$.

For MAR, these values were tightly clustered between 1.768 and 1.830. The between-imputation variance was only $0.000493$, so the five plausible completed datasets led to similar regression estimates. For MCAR, the five values ranged from 1.608 to 1.724, giving a larger between-imputation variance of $0.001923$. For MNAR, the range was approximately 1.688 to 1.764, with between-imputation variance $0.001067$.

Rubin's rules then combine within-imputation and between-imputation uncertainty. The main conceptual point is that the uncertainty associated with not knowing the missing values is explicitly included in the inferential calculation.

\subsection{Why MNAR Needs Extra Caution}
The MNAR experiment deliberately violates the usual MAR-type ignorability assumption used to justify standard MICE/EM-style analyses. Here, the missingness probability depends on the original $Y$ value itself.

The methods can still produce estimates that look numerically reasonable. For example, MICE gives $1.722$. But this does not show that MICE has ``solved MNAR.'' Without an explicit model or sensitivity assumptions describing how the unobserved values affect missingness, MNAR is not identified by the observed data alone.

The course material makes this point directly: increasing the number of imputations or iterations cannot repair a violated missingness assumption. The issue is not simply computation; it is the underlying information that is unavailable once the value is missing.

\subsection{Point Estimate Versus Inference}
The confidence intervals demonstrate why point estimates alone are not enough. Under MAR, MICE produces the smallest absolute bias in this simulation, but its 95\% interval ends at 1.982 and therefore does not contain the true value 2. Stochastic regression gives a somewhat more biased point estimate of 1.759, but its interval extends to 2.065.

Thus, closeness to the truth and uncertainty around the estimate are different aspects of an analysis. A method can have a relatively close point estimate and still yield an interval that does not contain the true value in a particular experiment.

% ---------------- 9 ----------------
\section{Inference Summary}
The central inferential quantities are summarized below.

\begin{table}[H]
\centering
\small
\caption{Inference summary for $\beta_1$. CI = 95\% confidence interval.}
\begin{tabular}{lrrrr}
\toprule
Mechanism & Method & $\hat\beta_1$ & SE & 95\% CI\\
\midrule
MCAR & Complete Case & 1.785 & 0.136 & (1.527, 2.025)\\
MCAR & Mean & 1.089 & 0.229 & (0.649, 1.565)\\
MCAR & Hot-Deck & 1.658 & 0.165 & (1.262, 1.918)\\
MCAR & Regression & 1.785 & 0.136 & (1.527, 2.025)\\
MCAR & Stochastic Regression & 1.706 & 0.140 & (1.487, 2.046)\\
MCAR & MICE & 1.659 & 0.114 & (1.433, 1.885)\\
MCAR & EM & 1.785 & 0.136 & (1.527, 2.025)\\
\midrule
MAR & Complete Case & 1.777 & 0.136 & (1.532, 2.055)\\
MAR & Mean & 1.307 & 0.186 & (0.915, 1.659)\\
MAR & Hot-Deck & 1.787 & 0.148 & (1.416, 1.996)\\
MAR & Regression & 1.777 & 0.136 & (1.532, 2.055)\\
MAR & Stochastic Regression & 1.759 & 0.147 & (1.499, 2.065)\\
MAR & MICE & 1.797 & 0.094 & (1.612, 1.982)\\
MAR & EM & 1.777 & 0.136 & (1.532, 2.055)\\
\midrule
MNAR & Complete Case & 1.749 & 0.122 & (1.501, 1.981)\\
MNAR & Mean & 1.413 & 0.170 & (1.047, 1.702)\\
MNAR & Hot-Deck & 1.722 & 0.123 & (1.469, 1.954)\\
MNAR & Regression & 1.749 & 0.122 & (1.501, 1.981)\\
MNAR & Stochastic Regression & 1.719 & 0.124 & (1.517, 1.987)\\
MNAR & MICE & 1.722 & 0.108 & (1.510, 1.933)\\
MNAR & EM & 1.749 & 0.122 & (1.501, 1.981)\\
\bottomrule
\end{tabular}
\end{table}

This table is intended to be read vertically within each mechanism. The same missing-data condition is held fixed while the methods are compared.

% ---------------- 10 ----------------
\section{Limitations}
\subsection{Single Simulated Dataset}
Only one complete simulated dataset and one realized missingness pattern were used for each mechanism. The results therefore describe the behavior of the methods in this controlled experiment rather than their long-run performance.

\subsection{No Monte Carlo Coverage Study}
The project stops at the bootstrap stage. It does not repeatedly generate new complete datasets and calculate how often 95\% confidence intervals contain the true parameter. Therefore, no empirical coverage probability is reported.

\subsection{Simple Data-Generating Model}
The model is linear, with normally distributed predictors and errors. Real datasets can contain nonlinear relationships, interactions, correlated predictors, multiple incomplete variables and non-normal outcomes.

\subsection{Simplified MICE Setting}
Only $Y$ is incomplete. A full MICE problem often contains multiple incomplete variables that are imputed sequentially. The present implementation is therefore a simplified PMM-based demonstration of multiple imputation and Rubin's rules.

\subsection{MNAR Identification}
The MNAR mechanism cannot generally be identified from the observed data without additional assumptions. A full MNAR study would require sensitivity analysis or an explicit missingness model.

\subsection{Bootstrap Scope}
Bootstrap inference estimates uncertainty relative to the observed sample. It does not by itself solve identification problems caused by an incorrect missingness assumption.

% ---------------- 11 ----------------
\section{Conclusion}
This project compared seven missing-data handling approaches within each of three missingness mechanisms. The mechanism was treated as a condition of the experiment, while the handling method was treated as the factor being compared.

The simulation shows that the choice of method can substantially affect regression estimates and inference. Mean imputation produced the largest downward distortion of the focal $\beta_1$ coefficient in MCAR, MAR and MNAR. The other methods produced estimates that were more tightly clustered around the complete-data relationship in this particular experiment, although their uncertainty intervals were not identical.

MICE added an important inferential perspective by generating several plausible completed datasets and incorporating both within-imputation and between-imputation uncertainty through Rubin's rules. EM provided a likelihood-based alternative that, in this special one-incomplete-variable setting, produced the same coefficient estimates as complete-case and regression-imputation analyses.

The MNAR experiment also highlighted a fundamental limitation: an algorithm cannot compensate for a missingness assumption that does not hold. When missingness depends on the unobserved outcome itself, additional modeling assumptions or sensitivity analysis are needed.

The broader lesson is therefore:
\[
\boxed{\text{Missing-data handling is part of statistical inference, not merely data cleaning.}}
\]

% ---------------- APPENDIX ----------------
\newpage
\appendix
\section{Additional MICE Details}
The MICE results are particularly useful for understanding how multiple imputation works numerically.

\begin{table}[H]
\centering
\small
\caption{Five MICE $\beta_1$ estimates and pooled results.}
\begin{tabular}{lrrrrrrrr}
\toprule
Mechanism & $m_1$ & $m_2$ & $m_3$ & $m_4$ & $m_5$ & Pooled & $W$ & $B$\\
\midrule
MCAR & 1.679 & 1.645 & 1.608 & 1.641 & 1.724 & 1.659 & 0.010727 & 0.001923\\
MAR & 1.768 & 1.830 & 1.794 & 1.799 & 1.792 & 1.797 & 0.008294 & 0.000493\\
MNAR & 1.764 & 1.688 & 1.748 & 1.708 & 1.700 & 1.722 & 0.010295 & 0.001067\\
\bottomrule
\end{tabular}
\end{table}

For MCAR,
\[
W=0.010727,
\qquad
B=0.001923,
\qquad
SE_{pooled}=0.114170.
\]

For MAR,
\[
W=0.008294,
\qquad
B=0.000493,
\qquad
SE_{pooled}=0.094262.
\]

For MNAR,
\[
W=0.010295,
\qquad
B=0.001067,
\qquad
SE_{pooled}=0.107590.
\]

These numbers make the role of multiple imputation concrete. The within-imputation variance reflects ordinary regression uncertainty inside each completed dataset, while the between-imputation variance reflects the fact that different plausible imputations do not produce exactly the same regression coefficient.

\section{EM Convergence Summary}
The EM algorithm converged to the following $\beta_1$ estimates in the current implementation:

\begin{table}[H]
\centering
\caption{EM convergence summary.}
\begin{tabular}{lrr}
\toprule
Mechanism & Iterations to convergence & Final $\beta_1$\\
\midrule
MCAR & 20 & 1.784723\\
MAR & 20 & 1.776896\\
MNAR & 14 & 1.748686\\
\bottomrule
\end{tabular}
\end{table}

\begin{figure}[H]
\centering
\includegraphics[width=0.78\textwidth]{figures/em_convergence_beta1.png}
\caption{EM convergence of $\beta_1$ in the three missingness settings.}
\end{figure}

\section{Method Comparison at a Glance}
The following compact view emphasizes the central project logic.

\begin{table}[H]
\centering
\small
\begin{tabular}{p{0.19\textwidth}p{0.73\textwidth}}
\toprule
Setting & Question asked\\
\midrule
MCAR & Given this MCAR dataset, how do the seven methods differ?\\
MAR & Given this MAR dataset, how do the seven methods differ?\\
MNAR & Given this MNAR dataset, how do the seven methods differ?\\
\bottomrule
\end{tabular}
\end{table}

This framing is the basis for all tables, plots and interpretation in the report.

% ---------------- REFERENCES ----------------
\newpage
\section*{References}
\begin{enumerate}[leftmargin=*]
\item Dempster, A. P., Laird, N. M., and Rubin, D. B. (1977). Maximum likelihood from incomplete data via the EM algorithm. \textit{Journal of the Royal Statistical Society: Series B}, 39(1), 1--38.
\item Rubin, D. B. (1987). \textit{Multiple Imputation for Nonresponse in Surveys}. Wiley.
\item van Buuren, S. (2018). \textit{Flexible Imputation of Missing Data}, 2nd ed. Chapman \& Hall/CRC.
\item Koley, D. \textit{Missing Data -- Unit I: Rigorous Treatment}. Course material, Department of Statistics, St. Xavier's University, Kolkata.
\item Koley, D. \textit{Missing Data -- Unit II: Imputation Methods and Algorithmic Approaches}. Course material, Department of Statistics, St. Xavier's University, Kolkata.
\end{enumerate}

\end{document}